% moid_disc: slicing a disc into thin rings (Ch 14, fig:moiddisc)
% Disc of radius R seen obliquely (circles drawn as ellipses, y-radius = 0.32 x-radius),
% axis vertical through the centre.  The slice is a thin ring of radius r and width dr,
% mass dm = sigma 2 pi r dr; every point of it is at distance r from the axis, so it
% contributes r^2 dm.  Summing from 0 to R gives (1/2) M R^2.
% Radii are measured along the horizontal (the direction drawn without foreshortening).
% Shared look of the moid_* family: see moid_rod.tex.
\begin{tikzpicture}[x=1cm,y=1cm]
  \def\k{0.32}     % ellipse squash
  \def\R{2.5}      % disc radius
  \def\h{0.2}      % disc thickness
  \def\r{1.45}     % slice radius
  \def\dr{0.3}     % slice width (exaggerated)
  % axis below the disc (the part behind the disc is hidden)
  \draw[rotation axis] (0,-\h-\k*\R) -- (0,-\h-\k*\R-0.75);
  % the disc: rim band, then the top face
  \draw[body,fill=cBlock!35] (-\R,0) -- (-\R,-\h) arc[start angle=180,end angle=360,x radius=\R,y radius=\k*\R]
     -- (\R,0) arc[start angle=0,end angle=-180,x radius=\R,y radius=\k*\R] -- cycle;
  \draw[body,fill=cBlock!20] (0,0) ellipse[x radius=\R,y radius=\k*\R];
  % the ring slice
  \filldraw[slice,even odd rule] (0,0) ellipse[x radius=\r+\dr,y radius=\k*(\r+\dr)]
                                 (0,0) ellipse[x radius=\r,y radius=\k*\r];
  % axis above the disc
  \draw[rotation axis] (0,0) -- (0,1.55) node[anchor=south] {axis};
  \fill[cInk] (0,0) circle[radius=0.04];
  % r: from the axis to the ring, along the left
  \draw[dim] (0,0) -- node[anchor=south,cSteel,inner sep=2pt] {$r$} (-\r,0);
  % dr: the ring's width, on the right
  \draw[{smallhead}-,thin line,cAcc] (\r,0) -- ++(-0.4,0);
  \draw[{smallhead}-,thin line,cAcc] (\r+\dr,0) -- (\R+0.25,0)
      node[anchor=west,inner sep=1.5pt] {$dr$};
  % dm, with a leader to the front of the ring
  \draw[leader,cAcc] (-40:{\r+\dr/2} and {\k*(\r+\dr/2)}) -- (2.6,-0.85)
      node[anchor=west,inner sep=1.5pt] {$dm=\sigma\,2\pi r\,dr$};
  % R: below the disc, on the left (the right is kept for dm)
  \draw[dimtick] (-\R,-\h-0.06) -- (-\R,-\h-\k*\R-0.6);
  \draw[dim] (-\R,-\h-\k*\R-0.48) -- node[anchor=north,cSteel] {$R$} (0,-\h-\k*\R-0.48);
\end{tikzpicture}
